% Formalization of a Kubernetes-like scheduling objective based on weighted scores
\documentclass[11pt]{article}
\usepackage[utf8]{inputenc}
\usepackage[T2A]{fontenc}
\usepackage[russian,english]{babel}
\usepackage{amsmath,amssymb}
\usepackage{hyperref}

\title{Формализация Kubernetes-like Scheduling\\(Weighted Placement Score Under Constraints)}
\author{cloud\_server\_optimization}
\date{\today}

\begin{document}
\maketitle

\section*{1. Идея}

В отличие от классической задачи bin packing, default \texttt{kube-scheduler}
обычно не решает задачу вида
\[
  \min \sum_n y_n.
\]
Его логика ближе к схеме:
\begin{enumerate}
  \item отфильтровать недопустимые узлы;
  \item посчитать score для каждого допустимого узла;
  \item выбрать узел с максимальным итоговым weighted score.
\end{enumerate}

Поэтому естественная математическая постановка имеет вид
\[
  \max \text{weighted placement score}
  \qquad \text{при ограничениях ресурсов и policy constraints.}
\]

\section*{2. Обозначения}

\paragraph{Поды.}
Пусть \(\mathcal{P}=\{1,\dots,P\}\) --- множество pod-ов, которые нужно разместить.
Для каждого \(p\in\mathcal{P}\):
\begin{align*}
  c_p &> 0 && \text{CPU request,}\\
  m_p &> 0 && \text{memory request.}
\end{align*}

\paragraph{Узлы.}
Пусть \(\mathcal{N}=\{1,\dots,N\}\) --- множество доступных узлов.
Для каждого \(n\in\mathcal{N}\):
\begin{align*}
  C_n &> 0 && \text{доступный CPU,}\\
  M_n &> 0 && \text{доступная память,}\\
  P_n &\in \mathbb{Z}_{>0} && \text{максимум pod-ов на узле.}
\end{align*}

\paragraph{Score plugins.}
Пусть \(\mathcal{K}=\{1,\dots,K\}\) --- множество scoring plugins
(например, \texttt{NodeResourcesFit}, \texttt{NodeResourcesBalancedAllocation},
\texttt{ImageLocality}, \texttt{PodTopologySpread}).

Для каждого plugin \(k\in\mathcal{K}\) задан вес \(w_k \ge 0\).
Пусть \(s_{p,n}^{(k)}\) --- score, который plugin \(k\) присваивает размещению
pod-а \(p\) на узле \(n\).
Обычно score нормируется, например, в диапазон \([0,100]\).

\paragraph{Feasibility indicator.}
Пусть
\[
  a_{p,n}\in\{0,1\},
\]
где \(a_{p,n}=1\), если узел \(n\) проходит фазу filtering для pod-а \(p\),
и \(a_{p,n}=0\) иначе.

Это место, где в модель попадают hard constraints:
resource requests, taints/tolerations, node affinity, volume constraints,
topology rules и другие требования.

\section*{3. Переменные}

Введём бинарные переменные размещения:
\[
  x_{p,n}\in\{0,1\},
\]
где \(x_{p,n}=1\), если pod \(p\) размещён на узле \(n\).

\section*{4. Итоговый weighted score}

Для пары \((p,n)\) введём post-placement величины:
\begin{align*}
  \hat{u}^{cpu}_{p,n} &= \frac{U_n^{cpu} + c_p}{C_n}, \\
  \hat{u}^{mem}_{p,n} &= \frac{U_n^{mem} + m_p}{M_n},
\end{align*}
где \(U_n^{cpu}\) и \(U_n^{mem}\) --- уже занятые ресурсы на узле \(n\) до
размещения pod-а \(p\).

\paragraph{NodeResourcesFit cо стратегией LeastAllocated.}
Чем больше свободных ресурсов останется после размещения, тем выше score:
\[
  s_{p,n}^{\mathrm{fit}} =
  \frac{100}{2}
  \left(
    \frac{C_n - U_n^{cpu} - c_p}{C_n}
    +
    \frac{M_n - U_n^{mem} - m_p}{M_n}
  \right).
\]

\paragraph{NodeResourcesBalancedAllocation.}
Этот plugin поощряет близкие доли загрузки CPU и памяти:
\[
  s_{p,n}^{\mathrm{bal}} =
  100\left(1 - \left|\hat{u}^{cpu}_{p,n} - \hat{u}^{mem}_{p,n}\right|\right).
\]

\paragraph{ImageLocality.}
Пусть
\[
  I_{p,n} \in \{0,1\}
\]
и \(I_{p,n}=1\), если на узле \(n\) уже есть pod того же image/service class,
что и у pod-а \(p\). Тогда:
\[
  s_{p,n}^{\mathrm{img}} = 100 \cdot I_{p,n}.
\]

\paragraph{Суммарный score.}
Для пары \((p,n)\) определим суммарный score:
\[
  S_{p,n}
  =
  w_{\mathrm{fit}} s_{p,n}^{\mathrm{fit}}
  +
  w_{\mathrm{bal}} s_{p,n}^{\mathrm{bal}}
  +
  w_{\mathrm{img}} s_{p,n}^{\mathrm{img}}.
\]

Тогда естественная целевая функция:
\[
  \max \sum_{p\in\mathcal{P}} \sum_{n\in\mathcal{N}} S_{p,n}\, x_{p,n}.
\]

Именно это и означает запись
\[
  \text{maximize weighted placement score under constraints.}
\]

\section*{5. Ограничения}

\paragraph{Каждый pod размещается не более чем на одном узле:}
\[
  \forall p\in\mathcal{P}:\quad
  \sum_{n\in\mathcal{N}} x_{p,n} \le 1.
\]

Если предполагается, что каждый pod обязан быть размещён, то можно писать:
\[
  \forall p\in\mathcal{P}:\quad
  \sum_{n\in\mathcal{N}} x_{p,n} = 1.
\]

\paragraph{Размещать pod можно только на допустимом узле:}
\[
  \forall p,n:\quad x_{p,n} \le a_{p,n}.
\]

\paragraph{Ограничения CPU, памяти и числа pod-ов на узле:}
\begin{align*}
  \forall n\in\mathcal{N}:\quad
  &\sum_{p\in\mathcal{P}} c_p\, x_{p,n} \le C_n, \\
  \forall n\in\mathcal{N}:\quad
  &\sum_{p\in\mathcal{P}} m_p\, x_{p,n} \le M_n, \\
  \forall n\in\mathcal{N}:\quad
  &\sum_{p\in\mathcal{P}} x_{p,n} \le P_n.
\end{align*}

\paragraph{Дополнительные policy constraints.}
При необходимости можно добавить ограничения общего вида:
\[
  g_j(x) \le b_j,
  \qquad j=1,\dots,J,
\]
где \(g_j(x)\) кодируют:
\begin{itemize}
  \item pod anti-affinity,
  \item topology spread constraints,
  \item размещение только на узлах с нужными labels,
  \item запрет на размещение на unschedulable nodes,
  \item budget на миграции или preemption.
\end{itemize}

\section*{6. Полная статическая постановка}

В итоге получается задача:
\[
  \max_{x} \sum_{p\in\mathcal{P}} \sum_{n\in\mathcal{N}} S_{p,n}\, x_{p,n}
\]
при ограничениях
\begin{align*}
  \forall p:\quad &\sum_n x_{p,n} \le 1, \\
  \forall p,n:\quad &x_{p,n} \le a_{p,n}, \\
  \forall n:\quad &\sum_p c_p\,x_{p,n} \le C_n, \\
  \forall n:\quad &\sum_p m_p\,x_{p,n} \le M_n, \\
  \forall n:\quad &\sum_p x_{p,n} \le P_n, \\
  \forall j:\quad &g_j(x) \le b_j, \\
  \forall p,n:\quad &x_{p,n}\in\{0,1\}.
\end{align*}

\section*{7. Локальная постановка для одного pod}

Если нужно записать именно поведение default scheduler для одного incoming pod
\(p^\star\), то задача ещё проще.

Введём переменные
\[
  x_n \in \{0,1\},
\]
где \(x_n=1\), если pod \(p^\star\) отправляется на узел \(n\).

Пусть
\[
  S_n = \sum_{k\in\mathcal{K}} w_k\, s_n^{(k)}
\]
--- итоговый score узла \(n\) для этого pod-а.

Тогда:
\[
  \max \sum_{n\in\mathcal{N}} S_n\, x_n
\]
при ограничениях
\begin{align*}
  &\sum_{n\in\mathcal{N}} x_n = 1, \\
  &x_n \le a_n \qquad \forall n, \\
  &c_{p^\star} x_n \le C_n^{\text{free}} \qquad \forall n, \\
  &m_{p^\star} x_n \le M_n^{\text{free}} \qquad \forall n, \\
  &x_n \in \{0,1\}.
\end{align*}

Это и есть самая близкая к Kubernetes запись:
\begin{quote}
выбрать допустимую ноду с максимальным weighted score.
\end{quote}

\section*{8. Связь с вашим проектом}

Если вы захотите приблизить текущую модель к \texttt{kube-scheduler}, то вместо
цели
\[
  \min \sum_n y_n
\]
можно перейти к формуле
\[
  \max \sum_{p,n} x_{p,n}
  \left(
    w_{\text{fit}} s_{p,n}^{\text{fit}}
    +
    w_{\text{bal}} s_{p,n}^{\text{bal}}
    +
    w_{\text{img}} s_{p,n}^{\text{img}}
  \right),
\]
где:
\begin{align*}
  s_{p,n}^{\text{fit}} &\text{ соответствует NodeResourcesFit(LeastAllocated),} \\
  s_{p,n}^{\text{bal}} &\text{ соответствует NodeResourcesBalancedAllocation,} \\
  s_{p,n}^{\text{img}} &\text{ соответствует ImageLocality.}
\end{align*}

Тогда ваш scheduler перестанет быть pure bin packing и станет ближе к
Kubernetes-style score-based placement.

\paragraph{Весовые профили стратегий.}
Если обозначить через \(\mathbf{w}^{(r)}\), \(\mathbf{w}^{(f)}\),
\(\mathbf{w}^{(h)}\) веса для стратегий Reactive, FFD и Hybrid, то можно писать:
\begin{align*}
  S_{p,n}^{(r)}
  &=
  4\,s_{p,n}^{\text{fit}}
  +
  1\,s_{p,n}^{\text{bal}}
  +
  1\,s_{p,n}^{\text{img}}, \\
  S_{p,n}^{(f)}
  &=
  2\,s_{p,n}^{\text{fit}}
  +
  3\,s_{p,n}^{\text{bal}}
  +
  1\,s_{p,n}^{\text{img}}, \\
  S_{p,n}^{(h)}
  &=
  3\,s_{p,n}^{\text{fit}}
  +
  2\,s_{p,n}^{\text{bal}}
  +
  2\,s_{p,n}^{\text{img}}.
\end{align*}

Тогда:
\begin{itemize}
  \item \textbf{Reactive} сохраняет порядок прихода pod-ов и выбирает узел по \(S_{p,n}^{(r)}\);
  \item \textbf{FFD} сначала сортирует крупные pod-ы и затем выбирает узел по \(S_{p,n}^{(f)}\);
  \item \textbf{Hybrid} также сортирует крупные pod-ы, но сильнее учитывает image locality через \(S_{p,n}^{(h)}\).
\end{itemize}

\end{document}
